\section{Messung von Impulsantworten} \label{ir_msrmt}
Um die akustischen Eigenschaften eines physischen Raumes in Form einer Impulsantwort abbilden zu können, wird der Raum
zunächst durch einen Messton angeregt. Dieser muss das gesamte Frequenzspektrum, auf das die IR später angewandt werden soll, enthalten \cite{smpm2001}.
Der im Raum resultierende Raumschall (der aus den Eigenschaften des Raumes sowie der Beschaffenheit
des anregenden Tons entstanden ist) wird aufgenommen. Das aufgenommene Signal resultiert aus der Transferfunktion des gemessenen Raumes,
angewandt auf den Messton. Da Raumhall-Effekte sich wie in \ref{irl_lti} dargestellt durch ihre Impulsantwort annähern lassen, kann die IR des Raumes
mithilfe entsprechender Verfahren aus diesem Signal berechnet werden \cite{novak2016}. Dabei macht es für die am Ende des Entfaltungsprozesses resultierende Impulsantwort
einen entscheidenden Unterschied, wie das Signal, mit dem der Raum angegeregt wird, beschaffen ist. In der Vergangenheit wurden oft sogenannte
pulsive Schallquellen genutzt \cite{farina2007b}, etwa Schüsse oder platzende Luftballons. Mit der steigenden Verbreitung digitaler Systeme
haben sich Mitte der 1980er Jahre zwei Schallquellen, die Maximum-Length-Sequence (kurz \textit{MLS}) und die Time Delay Spectrometry (kurz \textit{TDS}),
bei der ein linearer Sinus-Sweep über das zu messende Frequenzspektrum zum Einsatz kommt, durchgesetzt. Mit diesen Methoden konnten
verbesserter Rauschabstand und flachere Frequenzgänge der Impulsantworten erzielt werden \cite{farina2007b}. Seit Anfang der 2000er wird vermehrt
an Verfahren mit Einsatz von exponentiellen Sinus-Sweeps (kurz \textit{ESS}) und Logarithmischen Sinus-Sweeps (kurz \textit{LSS})
geforscht. Diese Verfahren ermöglichen verbesserte Abbildung der gemessenen Systeme besonders bei Vorhandensein von Nonlinearen oder 
Zeit-Varianten Eigenschaften \cite{farina2007b}.

\paragraph{Time Delay Spectrometry}
Bei TDS-basierten Messverfahren wird ein Sinus-Sweep genutzt, bei dem die Frequenz eines Sinus-Messsignals linear durch den zu messenden
Frequenzbereich, über einen bestimmten Zeitraum (die Messdauer), verändert wird. TDS-Verfahren sind allerdings aufgrund einiger Limitationen
in der praktischen Umsetzung besser für die Messung kurzer Impulsantworten geeignet, etwa denen von Lautsprecher-Systemen. Für die akkurate
Messung von längeren Impulsantworten müssten Impulse von unpraktikabler Länge, oder entsprechende Ausgleichsverfahren, auf die
wie hier nicht genauer eingehen können, genutzt werden \cite{farina2000, holters2009}. 
\paragraph{Minimum Length Sequence}
MLS sind periodische Sequenzen von binären Ziffern mit der Länge $L=2^N-1$. Bei der Nutzung einer MLS als Messimpuls wird der Raum durch den bei
der DA- und Schallwandlung erzeugten Schallimpuls angeregt. Aufgrund eines niedriger Energie in kurzen Impulsen dieser Art müssen in der Regel viele Messungen
vorgenommen und das Ergebnise gemittelt werden. Unter Anwendung entsprechender Prozesse weist das Ergebnis einer MLS-Messung bessere Eigenschaften
in der Immunität gegenüber Rauschen auf als TDS-Messungen \cite{rife1989}. Allerdings sind MLS-basierte Messverfahren sehr anfällig gegenüber 
Nonlinearitäten im gemessenen System und benötigen somit besondere Sorgfalt beim Einstellen der Messpegel \cite{holters2009}.
\paragraph{Exponential/Logarithmic Sine Sweep}
TDS- und MLS-Verfahren zur Messung der IR eines Systems haben beide nachteilhafte Eigenschaften, wenn in den gemessenen Systemen nonlineare Eigenschaften vorliegen.
Diese äußern sich in Form von Zeitaliasing-Artefakten, selbst wenn die Länge des Impulses angemessen gewählt wurde. Diese sind in der entfalteten
IR als unerwartete Peaks zu erkennen. Messverfahren mit linearen Sinus-Sweeps wie die TDS-Methode verbergen diese Peaks. Sie äußern sich bei diesen Verfahren in Form
eines mit dem Eingangssignal korrellierten Rauschens, welches nicht durch wiederholte Messung und Mittelwertbildung reduziert werden kann. Werden stattdessen
exponentielle Sinus-Sweeps, bei denen die momentäre Frequenz sich in Abhängigkeit der Zeit exponentiell verändert genutzt, dann sind diese Artefakte in der IR wieder gut
erkennbar \cite{farina2000}. Angelo Farina beschreibt ein Entfaltungs-Verfahren, mit dem diese Artefakte einfach aus der IR herausgerechnet werden können.
Unter Anwendung dieser Methodik ist sind ESS/LSS-Verfahren die geeignetsten Messverfahren für die Messung der Impulsantworten von Räumen.

\section{Entfaltung (Deconvolution) von Messungen} \label{deconvolution}
Hat man mit einem der in \ref{ir_msrmt} beschriebenen Verfahren das Ausgangssignal eines Raums gemessen, dann ist der nächste wichtige Schritt die Berechnung der
Impulsantwort des Raumes. Laut dem Faltungssatz ist eine Convolution in der Zeitdomäne äquivalent zu einer Multiplikation in der Frequenzdomäne \cite{oppenheim1999}.
Naiv liegt es damit nahe, dass eine Deconvolution in der Zeitdomäne äquivalent zu einer Division in der Frequenzdomäne ist. Mathematisch ist diese Annahme auch korrekt,
allerdings treten in der praktischen Anwendung gewisse Probleme auf. 

